\section{Data and calibration}\label{sec:q:data}

This appendix is what makes the calibrated set of Table~\ref{tab:q:data:parameters} reproducible: one subsection per block of the identification table of Section~\ref{sec:q:cal:blocks}, each stating the source, what is taken from it, the transformation in words, and the resulting number in a generated table. The transformations themselves are the pipeline scripts under \texttt{data/build/}, one per block, and provenance (citation, URL, access date, licence) is in \texttt{data/sources/}, which the tables cite rather than repeat. Every parameter is written to \texttt{data/processed/calibration.json} with a point, a range where the evidence supports one, and the subsection that identifies it; the model reads that file and nothing else. Units throughout are those of Section~\ref{sec:q:model:scope}: goods in trillions of 2011 international dollars, material in gigatonnes, one period one year, and pollution in gigatonnes of material. The base year is 1900, the first year the material flow accounts cover on one consistent accounting.

%% GENERATED by data/build/appendix_tables.py on 2026-09-17
%% Do not edit by hand: the file is rewritten by the next run.
\begin{table}[!htb]
\centering
\begin{threeparttable}
\caption{The calibrated parameter set.}\label{tab:q:data:parameters}
\begin{tabular}{llrrr}
\toprule
Symbol & Identifier & Value & Low & High \\
\midrule
$\Delta$ & \texttt{dt} & 1 & -- & -- \\
$\rho$ & \texttt{rho} & 0.03744 & -- & -- \\
$\eta$ & \texttt{eta} & 1.409 & 1.409 & 1.409 \\
$\psi$ & \texttt{psi\_v} & 0 & -- & -- \\
$A_0$ & \texttt{A0} & 1.365 & -- & -- \\
$g_A$ & \texttt{gA} & 0.01736 & 0.01736 & 0.02332 \\
$B_0$ & \texttt{B0} & 1 & -- & -- \\
$g_B$ & \texttt{gB} & 0 & -- & -- \\
$\beta$ & \texttt{beta\_s} & 0.7555 & -- & -- \\
$\mu_F$ & \texttt{mu\_F} & 0.4133 & 0.3761 & 0.4578 \\
$\varsigma$ & \texttt{sigma\_s} & 1 & -- & -- \\
$\kappa$ & \texttt{kappa} & $6.777\times10^{-6}$ & $2.391\times10^{-6}$ & $1.370\times10^{-5}$ \\
$\delta$ & \texttt{delta} & 0.03753 & 0.03148 & 0.04649 \\
$\bar R$ & \texttt{Rbar} & 0 & -- & -- \\
$\bar a$ & \texttt{abar} & 1 & -- & -- \\
$\xi$ & \texttt{xi} & 2.722 & 0.9073 & 8.166 \\
tail & \texttt{tail} & exp & -- & -- \\
$\psi$ (power tail) & \texttt{psi\_a} & 1 & -- & -- \\
$c_{N,0}$ & \texttt{cN0} & 0.01934 & -- & -- \\
$c_{N,\infty}$ & \texttt{cN\_inf} & 0.01406 & 0.007031 & 0.01406 \\
$g_{c_N}$ & \texttt{gcN} & 0.0317 & -- & -- \\
$\mu_N$ & \texttt{mu\_N} & 1.844 & 0 & 2.272 \\
$\kappa_N$ & \texttt{kap\_N} & 0.7281 & 0 & 7.281 \\
$\chi_N$ & \texttt{chi\_N} & 0.4818 & 0.2409 & 1 \\
$S_{\mathrm{ref}}$ & \texttt{Sref} & 13,429 & -- & -- \\
$c_{D,0}$ & \texttt{cD0} & $3.185\times10^{-4}$ & -- & -- \\
$c_{D,\infty}$ & \texttt{cD\_inf} & $3.185\times10^{-4}$ & -- & -- \\
$g_{c_D}$ & \texttt{gcD} & 0 & -- & -- \\
$\mu_D$ & \texttt{mu\_D} & 6.909 & 5.051 & 8.764 \\
$\kappa_D$ & \texttt{kap\_D} & 0.7281 & 0 & 7.281 \\
$\chi_D$ & \texttt{chi\_D} & 1 & -- & -- \\
$\bar X_{\max}$ & \texttt{Xmax} & 89,195 & 17,217 & $1.077\times10^{5}$ \\
$X_{\mathrm{ref}}$ & \texttt{Xref} & 13,503 & -- & -- \\
$\mu$ & \texttt{mu\_h} & 0.03489 & 0.03489 & 0.04596 \\
$c^c_0$ & \texttt{cc0} & 0.035 & 0.02 & 0.05 \\
$c^c_\infty$ & \texttt{cc\_inf} & 0.035 & -- & -- \\
$g_{c^c}$ & \texttt{gcc} & 0 & -- & -- \\
$c_{T,0}$ & \texttt{cT0} & 0.07 & 0.04 & 0.1 \\
$c_{T,\infty}$ & \texttt{cT\_inf} & 0.07 & -- & -- \\
$g_{c_T}$ & \texttt{gcT} & 0 & -- & -- \\
$\chi_T$ & \texttt{chi\_T} & 1.75 & 0.5 & 3 \\
$d^W$ & \texttt{dW} & 0.15 & 0.05 & 0.3 \\
$\theta_0$ & \texttt{theta0} & 0.01816 & 0.006523 & 0.01816 \\
$\theta_{\min}$ & \texttt{theta\_min} & 0.01816 & 0.006523 & 0.01816 \\
$\theta_P$ & \texttt{thetaP} & 0 & -- & -- \\
$\omega^Y$ & \texttt{omY} & 1 & -- & -- \\
$\omega^N$ & \texttt{omN} & 1 & -- & -- \\
$\omega^D$ & \texttt{omD} & 1 & -- & -- \\
$\omega^W$ & \texttt{omW} & 1 & -- & -- \\
$\omega^C$ & \texttt{omC} & 1 & -- & -- \\
$\phi^I$ & \texttt{phiI} & 4.611 & 2.043 & 6.907 \\
$\Omega^{N,S}$ & \texttt{OmNS} & 0.9322 & 0.7639 & 1.596 \\
$\Omega^{D,S}$ & \texttt{OmDS} & 0.009322 & 0 & 0.09322 \\
$\phi^W$ & \texttt{phiW} & 1 & -- & -- \\
$\phi^z$ & \texttt{phiz} & 1 & -- & -- \\
$\phi^P$ & \texttt{phiP} & 1 & -- & -- \\
$\phi^X$ & \texttt{phiX} & 1 & -- & -- \\
\bottomrule
\end{tabular}
\begin{tablenotes}[flushleft]
\footnotesize
\item[] \textit{Note:} The contents of \texttt{data/processed/calibration.json}; identifiers are those of \texttt{model/SYMBOLS.md}. A dash in the range columns is a point. Units: goods in trillion 2011 international dollars, material in gigatonnes, rates per year; cost scales in trillion dollars per Gt$^{1+\chi}$; $\xi$ in gigatonnes per trillion dollars; $\kappa$ per gigatonne of $P$.
\end{tablenotes}
\end{threeparttable}
\end{table}


\subsection{Accounting}\label{sec:q:data:accounting}

The accounting block is the best identified, because its objects are tonnes and the material flow accounts report tonnes. The source is the century-long global account of \textcite{krausmann2018} with the closed ledger of \textcite{haas2020}: extraction, secondary input, gross additions to stock and outflows by material category, annually from 1900 to 2015.

The durables coefficient $\phi^I$ is identified by the stored share $\sigma$ of Section~\ref{sec:q:model:within}: the model stores $\Omega\phi^IG$ of each year's material input in durables, whose observed counterpart is gross additions to stock over total material input. We set $\phi^I$ so that the century's cumulative gross additions to stock are reproduced at the model's activity aggregate, with the macro aggregates of Section~\ref{sec:q:data:production} and the pre-1950 use split held at the perpetual-inventory investment share. Table~\ref{tab:q:data:accounting} shows the stored share rising over the century, which a constant coefficient cannot follow: the coefficient read off the 1900 flows and the one read off the 2015 flows are the range, and the model starts with more of its 1900 input stored than the data show.

The relative intensities $\omega^j$ are not identified: no source reports waste generation by the use that generates it. \textcite{krausmann2018} report the outflow by type, and Table~\ref{tab:q:data:accounting} shows that split, which maps onto uses only coarsely, since the largest type, emissions, belongs to production and consumption alike. We hold the five intensities at a common value of one. Only the products $\Omega\omega^j$ enter the wedges of Section~\ref{sec:q:model:policy}, and Section~\ref{sec:q:cal:blocks} records why the split matters less than it looks; the perturbation that verifies it belongs to the review of the results.

Unused extraction has no counterpart in the material flow accounts, which count used extraction only. The ratio of hidden flows to commodities that the \textcite{eea2001} report for EU-15 domestic extraction by category, in the total material requirement accounting of the Wuppertal Institute, is the one published measurement obtained; $\Omega^{N,S}$ is that ratio weighted by the world composition of used extraction over 2000 to 2015. The range excludes biomass erosion at the low end and, at the high end, puts metals at the ratio the same source reports for the metals the EU imports, which is where the world's ore mining is. $\Omega^{D,S}$, mass displaced per tonne discovered, has no source: exploration moves drill core and trenches rather than overburden, and we set it at one hundredth of $\Omega^{N,S}$ with a range from zero to a tenth.

%% GENERATED by data/build/appendix_tables.py on 2026-09-17
%% Do not edit by hand: the file is rewritten by the next run.
\begin{table}[!htb]
\centering
\begin{threeparttable}
\caption{The accounting block.}\label{tab:q:data:accounting}
\begin{tabular}{lrl}
\toprule
Quantity & Value & Unit \\
\midrule
\multicolumn{3}{l}{\emph{The stored share $\sigma = \text{NAS}_{\text{gross}}/R$}} \\
1900 & 0.1639 & share \\
1950 & 0.2576 & share \\
2000 & 0.4776 & share \\
2015 & 0.5194 & share \\
cumulative 1900--2015 & 0.3992 & share \\
\multicolumn{3}{l}{\emph{The durables coefficient}} \\
$\sum\mathcal D/\sum G$, 1900--2015 & 6.939 & ratio \\
$\phi^I$ from the 1900 flows & 2.043 & ratio \\
$\phi^I$ from the 2015 flows & 6.907 & ratio \\
$\phi^I$, cumulative 1900--2015 (the point) & 4.611 & ratio \\
\multicolumn{3}{l}{\emph{The outflow by type in 2015, Krausmann et al.\ (2018)}} \\
processing waste & 9.447 & Gt/yr \\
emissions & 14.53 & Gt/yr \\
dissipative use & 6.124 & Gt/yr \\
excrements & 7.677 & Gt/yr \\
end-of-life waste & 10.52 & Gt/yr \\
water vapour & 9.319 & Gt/yr \\
$\omega^C/\omega^Y$, emissions unattributed & 2.003 & ratio \\
\multicolumn{3}{l}{\emph{Unused extraction per tonne used, EEA (2001) Table 3, and the 2000--2015 world shares of $N$}} \\
fossil fuels & 3.44 & t/t \\
\quad share of $N$ & 0.1732 & share \\
metal ores & 0.94 & t/t \\
\quad share of $N$ & 0.06385 & share \\
non-metallic minerals & 0.22 & t/t \\
\quad share of $N$ & 0.4914 & share \\
biomass & 0.62 & t/t \\
\quad share of $N$ & 0.2715 & share \\
EU-15 domestic total & 0.92 & t/t \\
$\Omega^{N,S}$ at the world composition (the point) & 0.9322 & t/t \\
$\Omega^{N,S}$, biomass erosion excluded (low) & 0.7639 & t/t \\
$\Omega^{N,S}$, metals at the EEA total ratio (high) & 1.596 & t/t \\
$\Omega^{D,S}$ (judgement) & 0.009322 & t/t \\
\bottomrule
\end{tabular}
\begin{tablenotes}[flushleft]
\footnotesize
\item[] \textit{Note:} Flows from Haas et al.\ (2020) and Krausmann et al.\ (2018), block B1; $\mathcal D$ at the common intensity $\omega^j=1$ and the macro aggregates of block B2 with the pre-1950 use split at the perpetual-inventory investment share. The EEA ratios are hidden flows per tonne of commodity for EU-15 domestic extraction in 1995; provenance in \texttt{data/sources/}.
\end{tablenotes}
\end{threeparttable}
\end{table}


\subsection{Production and trends}\label{sec:q:data:production}

The macro aggregates are those of block B2: world GDP from the Maddison Project Database \parencite{bolt2024}, the capital stock, use split, depreciation and labour share from the Penn World Table \parencite{feenstra2015}, with the capital stock before 1950 a perpetual inventory anchored at 1950 and flagged reconstructed. Before 1950 the split of output between consumption and gross formation uses the investment share the perpetual inventory needs to reproduce the 1950 anchor, so that the reconstructed capital stock and the reconstructed gross formation rest on one share.

The fixed factor $1-\mu_F$ stands in for labour, so $\mu_F$ is one minus the world labour share. The material elasticity $\gamma=\mu_F\left(1-\beta\right)$ is the value of primary material input over GDP in the two calibration years, 2011, the dollar base year, and 2015, the last year of the material block: the fossil masses of block B3 at the World Bank commodity prices \parencite{worldbank2026cmo}, metal ores and non-metallic minerals at the USGS unit values \parencite{usgs2024ds140}, and biomass at the value added of agriculture, forestry and fishing from the World Development Indicators \parencite{worldbank2026wdi}, all over nominal world GDP; $\beta$ follows. The extraction cost share of GDP, which Section~\ref{sec:q:data:extraction} uses for the level of $c_N$, is the same value less the natural resource rents of the same source, which are defined as revenue less production cost. Table~\ref{tab:q:data:production} shows the components. The biomass component is a value added rather than a harvest value, and the metal-ore component prices every tonne of ore at the iron ore unit value; both are approximations the table makes visible.

We hold $\varsigma=1$. Under Cobb--Douglas the material value share is constant and the intensity $R/Y$ falls one for one with the real material price, and Table~\ref{tab:q:data:production} reports the century's elasticity of the intensity to the mass-weighted real price of \textcite{jacks2019} as the test: the intensity fell faster than the price rose, so the data lean towards substitution above unity, but the metals price fell over the same century, and one elasticity across the whole aggregate is not identified from two trends. The long-run price elasticity of material demand that the identification table names cannot be read off a century in which income growth dominates the price effect.

The trends are over-identified only in principle. Under $\varsigma=1$ the production function depends on $A$ and $B$ through the one constant $AB^{\gamma}$, so any split of a given composite $g_A+\gamma g_B$ produces the same path; what growth accounting identifies is the composite, output growth less the capital and material elasticities times the growth of capital and material input, over 1900 to 2015 as the point and over 1950 to 2015 as the range. The split is a convention: $g_B=0$ and $g_A$ the composite, with $B_0=1$. The test the identification table promises is then the balanced-path one. On a circular path material input is constant and $\Omega$ falls at the rate output grows; on a balanced dematerialization path the ratio of the two rates is $\left(\mu_N-1\right)/\mu_N$. Table~\ref{tab:q:data:production} shows that the century is on neither: material input grew throughout, and the intensity on the model's definition fell at a fraction of the rate output grew. The history is the transition, not a balanced path, and the trends are calibrated on it as such. $A_0$ is set so that the production function reproduces 1900 output at the 1900 capital stock and material input, neglecting the damage factor at the 1900 pollution stock, which Section~\ref{sec:q:data:damages} shows to be a twentieth of a percent.

%% GENERATED by data/build/appendix_tables.py on 2026-09-18
%% Do not edit by hand: the file is rewritten by the next run.
\begin{table}[!htb]
\centering
\begin{threeparttable}
\caption{Production and trends.}\label{tab:q:data:production}
\begin{tabular}{lrl}
\toprule
Quantity & Value & Unit \\
\midrule
\multicolumn{3}{l}{\emph{Depreciation and the labour share, PWT 10.01 world aggregates}} \\
$\delta$, mean 1950--2019 & 0.03753 & per year \\
$\delta$ in 1950 & 0.03148 & per year \\
$\delta$ in 2019 & 0.04649 & per year \\
labour share, mean 1950--2019 & 0.5867 & share \\
\quad lowest decade mean & 0.5422 & share \\
\quad highest decade mean & 0.6239 & share \\
$\mu_F$ & 0.4133 & exponent \\
\multicolumn{3}{l}{\emph{The value of primary material input, nominal}} \\
fossil fuels, 2015 & 2.537 & trillion current US\$ \\
metal ores, 2015 & 0.5247 & trillion current US\$ \\
non-metallic minerals, 2015 & 0.4251 & trillion current US\$ \\
biomass (agriculture value added), 2015 & 3.174 & trillion current US\$ \\
natural resource rents, 2015 & 1.324 & trillion current US\$ \\
world GDP, 2015 & 75.84 & trillion current US\$ \\
value share of GDP, 2011 & 0.1143 & share of GDP \\
value share of GDP, 2015 & 0.08784 & share of GDP \\
cost share of GDP, 2011 & 0.06602 & share of GDP \\
cost share of GDP, 2015 & 0.07038 & share of GDP \\
$\gamma=\mu_F(1-\beta)$, the material elasticity & 0.1011 & elasticity \\
$\beta$ & 0.7555 & share \\
$\beta_K=\mu_F\beta$ & 0.3123 & elasticity \\
\multicolumn{3}{l}{\emph{Growth rates, log per year}} \\
$Y$, 1900--2015 & 0.02974 & per year \\
$K$, 1900--2015 & 0.03253 & per year \\
$R$, 1900--2015 & 0.02201 & per year \\
$R/Y$, 1900--2015 & -0.007736 & per year \\
$\Omega$ on the model's definition, 1900--2015 & -0.008398 & per year \\
real material price, mass-weighted, 1900--2015 & 0.004541 & per year \\
real metals price, 1900--2015 & -0.003417 & per year \\
$\mathrm d\ln(R/Y)/\mathrm d\ln p$, 1900--2015 & -1.704 & elasticity \\
$\mathrm d\ln(R/Y)/\mathrm d\ln p$, 1950--2015 & -2.361 & elasticity \\
\multicolumn{3}{l}{\emph{The trends}} \\
$g_A+\gamma g_B$, 1900--2015 (the point) & 0.01736 & per year \\
$g_A+\gamma g_B$, 1950--2015 & 0.02332 & per year \\
implied circular growth rate $g$ & 0.02524 & per year \\
$g_Y/(-g_\Omega)$, 1900--2015 & 3.542 & ratio \\
$A_0$ at $B_0=1$ & 1.365 & level \\
\bottomrule
\end{tabular}
\begin{tablenotes}[flushleft]
\footnotesize
\item[] \textit{Note:} Fossil masses by fuel from block B3 at the World Bank Pink Sheet prices, oil at 7.33 barrels per tonne (Rogner et al.\ 2012, footnote 5), gas at the IPCC (2006) net calorific value; metal ores at the USGS iron ore unit value and non-metallic minerals at the mean of crushed stone and construction sand and gravel; agriculture, rents and GDP from the World Development Indicators. Prices are Jacks (2019) indices at the B3 weights. Cobb--Douglas predicts an intensity elasticity of $-1$; on a circular path $g_Y/(-g_\Omega)=1$.
\end{tablenotes}
\end{threeparttable}
\end{table}


\subsection{Extraction and exploration}\label{sec:q:data:extraction}

\smalltitle{The composition of the reserve.} The reserve and the discovery ceiling are measured for the exhaustible categories, fossil fuels and metal ores, by block B3: cumulative extraction from 1900 to the reserve vintage plus the reserves at that vintage, from \textcite{rogner2012} for fossil fuels and the \textcite{usgs2026mcs} for ores, on the gross-ore basis where the source reports it. The model draws its one reserve down by total extraction $N$, three quarters of which is biomass and construction minerals, whose stock effect is nil over the horizon. Those flows cannot be put into $N$ and kept out of $S$ inside a one-material parameter set, so the composition chosen is the one that makes the stock effect on the aggregate match the exhaustible categories' own relative depletion: $S_0$ is the exhaustible reserve divided by the exhaustible share of cumulative extraction over 1900 to 2015. The relative depletion of the one reserve then tracks that of fossil fuels and ores, and $\left(S_{\mathrm{ref}}/S\right)^{\mu_N}$ on the aggregate is what it would be on the exhaustible stock. What is lost is that the renewable three quarters of extraction inherit a stock effect they do not have, and that the level of $S$ is a bookkeeping quantity rather than a tonnage anyone measures. The alternative, the exhaustible reserve unscaled, would be exhausted by the century's own extraction and is not admissible. The USGS gross-ore commodities cover a fraction of the material flow accounts' metal ores, and their reserves and resources are scaled up by the ratio of cumulative extraction in the two sources. $X_0$ adds the extraction before 1900, which is known for fossil fuels only, so it is a lower bound. $\bar X_{\max}$ is the ultimately recoverable resource of the same sources scaled the same way, with the low, central and high readings of \textcite{rogner2012} as the range and the conventional-plus-coal variant recorded beside it; Table~\ref{tab:q:data:extraction} shows how far the central reading sits above the low one, which is the whole question of what the coal and unconventional gas resources are worth.

\smalltitle{The stock-effect elasticity.} $\mu_N$ cannot be separated from technical progress in extraction by a price history alone, for the reason \theory{Section}{Balanced dematerialization: rate matching} and Section~\ref{sec:q:cal:blocks} give. We resolve it by an assumption stated as such: $c_N\left(t\right)$ is held to absorb the flow-scale effect of the effort cost, so that the trend of the real material price over 1900 to 2015 is the stock effect, and $\mu_N$ is minus the log change of the mass-weighted real price of \textcite{jacks2019} over the log change of the composed reserve. The range is the same construction at the low and high reserve figures and, at the bottom, the metals-only price channel, on which the real price fell while the reserve was drawn down and which therefore gives a negative elasticity, floored at zero. The physical channel, the elasticity of energy per tonne to ore grade of \textcite{mudd2010}, is shown beside them; it is of the order of the metals reading and not of the fossil one.

\smalltitle{The effort cost.} For a cost $c_N\left(\kappa_NN+N^{1+\chi_N}/\left(1+\chi_N\right)\right)$ the ratio of production cost to value at the margin is $\left(\kappa_N+N/\left(1+\chi_N\right)\right)/\left(\kappa_N+N\right)$, so $\chi_N$ is the value share of Section~\ref{sec:q:data:production} over the cost share less one, an upper reading because the Hotelling rent is also in the value; the range runs down to half of it and up to the quadratic. $\kappa_N$ is a tenth of 1900 extraction, which is a judgement and not a reading, and $S_{\mathrm{ref}}=S_0$. The level in 2015 is the cost share times output, and the implied path of $c_N$ over the century follows from the price index, the reserve path and the flow term; Table~\ref{tab:q:data:extraction} shows that it falls by a factor of about three with no sign of flattening, so the convergent form is pinned at both ends, the 1900s mean and the 2006 to 2015 mean, and only the rate is fitted, with the fit's error in logs reported. The floor at the end-of-sample level is the conservative reading, no progress after the sample, and half of it is the range.

\smalltitle{Exploration.} The MinEx series of \textcite{schodde2023} give the average cost of a discovery in two periods and the room below the contained-metal resource of the \textcite{usgs2026mcs}; $\mu_D$ is the elasticity of the one to the other, at the low, central and high resource figures. It is large, and it rests on a count of deposits rather than a tonnage, which the table's row labels say. The level $c_D$ comes from the exploration margin: at replacement, discoveries equal to extraction, the marginal discovery cost equals the resource rent per tonne, which the World Development Indicators give. Table~\ref{tab:q:data:extraction} shows the total exploration cost that implies against the world's metals exploration spend, which is an order of magnitude smaller; oil and gas exploration, which is the larger part of the world's spend, is not in the comparison. $\chi_D=1$ by convention, $\kappa_D$ a tenth of 1900 extraction, and $X_{\mathrm{ref}}=X_0$. By construction of $S_0$ no discovery is needed before the reserve vintage, so history says nothing about the level of exploration, and the convexity condition (C3) of the sufficiency appendix fails on both cost functions: the finite choke and the stock elasticities above the flow curvature break joint convexity, which Section~\ref{sec:q:data:checks} reports rather than hides.

%% GENERATED by data/build/appendix_tables.py on 2026-09-18
%% Do not edit by hand: the file is rewritten by the next run.
\begin{table}[!htb]
\centering
\begin{threeparttable}
\caption{Extraction and exploration.}\label{tab:q:data:extraction}
\begin{tabular}{lrl}
\toprule
Quantity & Value & Unit \\
\midrule
\multicolumn{3}{l}{\emph{The composition of the reserve}} \\
exhaustible share of cumulative $N$, 1900--2015 & 0.2364 & share \\
\quad in 1900 & 0.1539 & share \\
\quad in 2015 & 0.2359 & share \\
metal-ore coverage scale, material flows over USGS gross ore & 2.26 & ratio \\
fossil reserve, low & 1,805 & Gt \\
fossil reserve, high & 3,080 & Gt \\
metal-ore reserve, gross basis, scaled & 732.1 & Gt \\
exhaustible reserve, midpoint & 3,175 & Gt \\
$S_0$ (the point) & 13,429 & Gt \\
$S_0$, low & 10,734 & Gt \\
$S_0$, high & 16,123 & Gt \\
extraction before 1900, fossil & 17.6 & Gt \\
$X_0$ & 13,503 & Gt \\
\multicolumn{3}{l}{\emph{The discovery ceiling}} \\
fossil URR, low & 1,880 & Gt \\
fossil URR, central & 18,874 & Gt \\
fossil URR, high & 23,235 & Gt \\
\quad conventional plus coal, central & 15,844 & Gt \\
metal-ore resources, gross basis, scaled & 2,212 & Gt \\
$\bar X_{\max}$ (the point) & 89,195 & Gt \\
$\bar X_{\max}$, low & 17,217 & Gt \\
$\bar X_{\max}$, high & $1.077\times10^{5}$ & Gt \\
$\bar X_{\max}$, conventional plus coal & 76,378 & Gt \\
\multicolumn{3}{l}{\emph{The stock-effect elasticity}} \\
mass-weighted real price index, 2015 & 168.6 & 1900 = 100 \\
metals real price index, 2015 & 67.5 & 1900 = 100 \\
$\mu_N$ at $S_0$ low & 1.416 & elasticity \\
$\mu_N$ at $S_0$ midpoint (the point) & 1.844 & elasticity \\
$\mu_N$ at $S_0$ high & 2.272 & elasticity \\
$\mu_N$, metals only & -1.531 & elasticity \\
energy--grade elasticity, gold (Mudd 2010) & -0.2848 & elasticity \\
\multicolumn{3}{l}{\emph{The effort cost}} \\
$\chi_N$ & 0.4818 & exponent \\
$\kappa_N$ & 0.7281 & Gt/yr \\
$C^N$ in 2015 & 7.309 & trillion \$ \\
marginal cost in 2015 & 117.4 & US\$/t \\
$c_{N,0}$ & 0.01934 & trillion \$/Gt\^{}(1+chi\_N) \\
$c_{N,\infty}$ & 0.01406 & trillion \$/Gt\^{}(1+chi\_N) \\
$g_{c_N}$ & 0.0317 & per year \\
fit RMSE, log points & 0.2749 & log points \\
\multicolumn{3}{l}{\emph{Exploration}} \\
cost per discovery, 1975--2005 & 65 & 2023 US\$ million \\
cost per discovery, 2011--2020 & 218 & 2023 US\$ million \\
cumulative contained metal to 1990 & 1.219 & Gt \\
cumulative contained metal to 2015 & 2.363 & Gt \\
$\mu_D$ at the low URR & 5.051 & elasticity \\
$\mu_D$ at the central URR (the point) & 6.909 & elasticity \\
$\mu_D$ at the high URR & 8.764 & elasticity \\
resource rent per tonne of $N$, 2011--2015 & 0.02854 & trillion \$/Gt \\
$\kappa_D$ & 0.7281 & Gt/yr \\
$c_{D,0}$ & $3.185\times10^{-4}$ & trillion \$/Gt\^{}(1+chi\_D) \\
$C^D$ at $D=N_{2015}$ & 1.279 & trillion \$ \\
world metals exploration spend, 2012 & 0.0421 & trillion 2023 US\$ \\
\bottomrule
\end{tabular}
\begin{tablenotes}[flushleft]
\footnotesize
\item[] \textit{Note:} Reserves, resources and the price indices from block B3; extraction by category from Krausmann et al.\ (2018). The model's $S_0$, $X_0$ and $\bar X_{\max}$ are the exhaustible figures divided by the exhaustible share, so that the relative depletion of the one reserve tracks that of fossil fuels and ores. $\mu_N$ is $-\Delta\ln p/\Delta\ln S$ over 1900--2015 with $c_N(t)$ absorbing the flow-scale effect; $\chi_N$ is the value share over the cost share less one. Costs in trillion 2011 international dollars and gigatonnes; $c_N$ and $c_D$ in trillion dollars per Gt$^{1+\chi}$.
\end{tablenotes}
\end{threeparttable}
\end{table}


\subsection{Waste handling and recycling}\label{sec:q:data:waste}

\smalltitle{The stockpile.} The waste stock is read broadly, as \theory{Section}{Aggregate material balances} allows: everything ever discarded to controlled or uncontrolled disposal and not recovered, cumulated from 1900 in \textcite{haas2020}, for fossil, metal and non-metallic mineral outflows. Biomass residues cycle ecologically and are not stockpiled; the table shows the residue excluded. The handling share $\mu$ is set so that $\mu\mathcal W$ in 2015 equals the handled flow, the outflow of those three categories to disposal in the same year; the reading that counts the same year's recovered flow as handled too is the range's upper end. The specification tension between the model's geometric draw on the stockpile and the age-dependent recovery in the data is recorded in \texttt{notes/data/waste\_stock.md}.

\smalltitle{The treated share.} $\varpi$ is a control, not a parameter, but the solved path can be checked against it. Table~\ref{tab:q:data:waste} gives two readings: recovery over recovery plus disposal on the material flow accounts, which is a floor because the accounts do not separate controlled from uncontrolled disposal, and the municipal treated share of \textcite{unep2024} and \textcite{kaza2018}, which is a measurement of three percent of the handled flow.

\smalltitle{Costs.} Every unit cost in block B4 is a municipal solid waste cost, and municipal waste is three percent of the flow the model handles, so the income-group ladder of \textcite{kaza2018} fixes the shape of the handling cost and not its level. The shape is the ratio of the treatment scale $c_T$, the marginal treatment cost at full treatment, to the collection charge $c^c$, the intercept of the ladder at a low treated share: the high-income controlled-to-sanitary landfill midpoint over the low-income collection and transfer midpoint, converted at one thousandth of a trillion dollars per gigatonne, with neither source stating its price base and no deflation applied. The level is fitted, not measured. With the ratio held, one common factor on both charges is chosen so that the planner path solved from 1900 reproduces the 2015 treated share of the material flow accounts, recovery over recovery plus disposal, by bracketing and bisection on the factor; Table~\ref{tab:q:data:waste} gives the factor, the share achieved and the solves taken, the municipal midpoints as the alternative at the municipal scale, and the ranges of $c^c$ and $c_T$ are the ladder's ranges at the fitted factor. The factor lies inside the ladder's own range, so the municipal charges were of the right order for the aggregate after all; what they could not do at their midpoint was reproduce the date at which treatment begins. The second reading of 2015 circularity, secondary input over the outflows from use, is reported beside its target and not fitted: one factor cannot hit both, and the model's recycled share at the fitted factor is more than twice the observed one, because the yield the recycling margin chooses in 2015 is well above the observed yield through which $\xi$ is set, and because the model's handled flow in 2015 is several times the disposal flow $\mu$ was set to reproduce, the model's stockpile carrying every outflow and the unused extraction where the accounts' carries disposal alone. Both are the tail rate's and the stockpile's problem, not the charges', and are left to the sensitivity analysis over $\xi$. The curvature $\chi_T$ is what four confounded points on the ladder support, a range from one half to three with the midpoint as the point, and the leakage $d^W$ of the treated stream's residue, which nothing measures, is a range around the composition of the treated municipal stream.

\smalltitle{The ceiling and the tail.} The soft ceiling $\bar a=1$ is the baseline and the hard ceiling is the metals-only mass-weighted end-of-life recycling rate of \textcite{unepirp2011}, which is essentially the iron and steel figure; on a mass basis across all materials it is conservative rather than tight, for the reason Section~\ref{sec:q:cal:aggregation} gives. The tail rate $\xi$ is set through a point on the yield function rather than a gradient, since no source gives the cost of recovery as a function of the rate achieved: the observed yield is recovery over the treated flow on the material flow accounts at the municipal treated share, the slope $a'\left(x\right)$ is the reciprocal of the capital outlay per tonne of added annual recovery that the \textcite{epa2024} cost for the United States, and the exponential form gives $\xi=a'/\left(\bar a-a\right)$. Table~\ref{tab:q:data:waste} shows the recycling capital the point implies against the capital stock. The range is an order of magnitude around the point, and the power tail with exponent one, the largest that keeps the yield concave, is the sensitivity.

%% GENERATED by data/build/appendix_tables.py on 2026-09-17
%% Do not edit by hand: the file is rewritten by the next run.
\begin{table}[!htb]
\centering
\begin{threeparttable}
\caption{Waste handling and recycling.}\label{tab:q:data:waste}
\begin{tabular}{lrl}
\toprule
Quantity & Value & Unit \\
\midrule
\multicolumn{3}{l}{\emph{The stockpile and the handling share}} \\
$\mathcal W$ in 2015, fossil, metals, minerals & 533.7 & Gt \\
\quad biomass residue, excluded & 783 & Gt \\
outflow to disposal, 2015 & 18.62 & Gt/yr \\
recovered flow, 2015 & 5.905 & Gt/yr \\
$\mu$ (the point) & 0.03489 & per year \\
$\mu$ with recovery in the handled flow (high) & 0.04596 & per year \\
residence time $1/\mu$ & 28.66 & years \\
\multicolumn{3}{l}{\emph{The treated share, a check on the solved path}} \\
recovery over recovery plus disposal, B1 & 0.2407 & share \\
municipal treated share, UNEP (2024) & 0.62 & share \\
municipal treated share, Kaza et al.\ (2018) & 0.67 & share \\
\multicolumn{3}{l}{\emph{Costs: the shape, from the municipal ladder}} \\
collection, low income, low & 20 & US\$/t \\
collection, low income, high & 50 & US\$/t \\
$c^c$ at the municipal scale (the alternative) & 0.035 & trillion \$/Gt \\
sanitary landfill, high income, low & 40 & US\$/t \\
sanitary landfill, high income, high & 100 & US\$/t \\
$c_T$ at the municipal scale (the alternative) & 0.07 & trillion \$/Gt \\
$c_T/c^c$, held & 2 & ratio \\
\multicolumn{3}{l}{\emph{Costs: the level, fitted to 2015 circularity}} \\
factor on the ladder (the level) & 1.242 & ratio to the municipal ladder \\
solves taken & 8 & solves \\
treated share 2015 on the solved path (target: the B1 reading above) & 0.2522 & share \\
recycled share 2015, $RR/W$, B1 (target) & 0.09588 & share \\
recycled share 2015 on the solved path (reported, not fitted) & 0.22 & share \\
$a\varpi$ in 2015 on the solved path & 0.1734 & share \\
$c^c$ (the point) & 0.04347 & trillion \$/Gt \\
$c_T$ (the point) & 0.08693 & trillion \$/Gt \\
$C^W$ in 2015 at the municipal treated share & 0.8692 & trillion \$ \\
\multicolumn{3}{l}{\emph{The ceiling and the tail}} \\
$\bar a_H$, metals end-of-life recycling rate & 0.708 & share \\
\quad low & 0.5269 & share \\
\quad high & 0.8892 & share \\
observed yield $a$, 2015 & 0.3883 & share \\
capital per added tonne of recovery, EPA (2024) & 0.6006 & trillion \$ per Gt/yr \\
$a'(x)$ implied & 1.665 & Gt per trillion \$ \\
$\xi$ (the point) & 2.722 & Gt per trillion \$ \\
$\xi$, low & 0.9073 & Gt per trillion \$ \\
$\xi$, high & 8.166 & Gt per trillion \$ \\
recycling capital per treated tonne $x$, 2015 & 0.1806 & trillion \$/Gt \\
$K^R$ implied, 2015 & 2.746 & trillion \$ \\
\multicolumn{3}{l}{\emph{The 1900 stockpile}} \\
$\mathcal W_0$ consistent with the 1900 disposal flow (the point) & 16.39 & Gt \\
$\mathcal W_0$, the 1900 residue alone (low) & 0.572 & Gt \\
\bottomrule
\end{tabular}
\begin{tablenotes}[flushleft]
\footnotesize
\item[] \textit{Note:} Stocks and flows from Haas et al.\ (2020), block B1, without biomass; treatment shares and costs from Kaza et al.\ (2018) and UNEP (2024), block B4, US dollars per tonne with the price base not stated and converted at 0.001 trillion dollars per gigatonne; the level of the charges is the factor on the ladder's midpoints that makes the planner path solved from 1900 at $T=200$ reproduce the 2015 treated share (\texttt{data/build/c4\_handling\_level.jl}), and the ranges of $c^c$ and $c_T$ are the ladder's at that factor; the metals recycling rate from UNEP IRP (2011) weighted by USGS mine production; the capital outlay from US EPA (2024). $\chi_T=1.75$ in $[0.5,3]$ and $d^W=0.15$ in $[0.05,0.30]$ are judgements of block B4 and carry no row.
\end{tablenotes}
\end{threeparttable}
\end{table}


\subsection{Preferences}\label{sec:q:data:preferences}

On a balanced path the Euler equation reads $r=\rho+\eta g$, and the consumption--wealth ratio is $r-g$ whatever the preferences, so the two moments of the identification table pin one combination and not two. What separates them is time variation: the decade means of the world real return on capital of the Penn World Table are regressed on the decade means of world output growth over 1950 to 2019, $\eta$ the slope and $\rho$ the intercept, with the fit reported in Table~\ref{tab:q:data:preferences}. The data value of $\eta$ passes the collapse ranking condition $\left(1-\delta\right)^{1-\eta}<1+\rho$ at the data $\rho$ and the depreciation rate of Section~\ref{sec:q:data:production}; the table gives the largest value that would pass, which is what the theory would have required had it failed. The pure rate of time preference the regression returns is high by the standards of the integrated-assessment literature, because the return it is fitted to is the gross return on productive capital rather than a risk-free rate; the implied real rate on the circular path is shown for comparison. The utility channel of pollution is off, $\psi=0$, for the reason Section~\ref{sec:q:data:damages} gives.

%% GENERATED by data/build/appendix_tables.py on 2026-09-18
%% Do not edit by hand: the file is rewritten by the next run.
\begin{table}[!htb]
\centering
\begin{threeparttable}
\caption{Preferences.}\label{tab:q:data:preferences}
\begin{tabular}{lrl}
\toprule
Quantity & Value & Unit \\
\midrule
\multicolumn{3}{l}{\emph{Decade means, world}} \\
real return on capital, 1950s & 0.1059 & per year \\
\quad output growth & 0.04492 & per year \\
real return, 1960s & 0.1097 & per year \\
\quad output growth & 0.04978 & per year \\
real return, 1970s & 0.092 & per year \\
\quad output growth & 0.04077 & per year \\
real return, 1980s & 0.07952 & per year \\
\quad output growth & 0.03032 & per year \\
real return, 1990s & 0.08453 & per year \\
\quad output growth & 0.02958 & per year \\
real return, 2000s & 0.08554 & per year \\
\quad output growth & 0.03987 & per year \\
real return, 2010s & 0.08765 & per year \\
\quad output growth & 0.03652 & per year \\
\multicolumn{3}{l}{\emph{The Euler regression $r=\rho+\eta g$}} \\
$\rho$, intercept & 0.03744 & per year \\
$\eta$, slope & 1.409 & elasticity \\
$R^2$ & 0.8253 & nan \\
largest $\eta$ passing the collapse ranking condition & 1.961 & elasticity \\
$\eta$ used & 1.409 & elasticity \\
$\rho$ used & 0.03744 & per year \\
implied real rate on the circular path & 0.075 & per year \\
consumption--wealth ratio $C/K$, 2015 & 0.1744 & per year \\
\bottomrule
\end{tabular}
\begin{tablenotes}[flushleft]
\footnotesize
\item[] \textit{Note:} The return is the PWT 10.01 internal rate of return \texttt{irr}, weighted by the capital stock at current PPPs; growth is the log change of Maddison world GDP. On a balanced path the consumption--wealth ratio equals $r-g$ for any preferences, so it is a check and not a second moment; the decade variation is what separates $\rho$ from $\eta$.
\end{tablenotes}
\end{threeparttable}
\end{table}


\subsection{Damages and decay}\label{sec:q:data:damages}

Under decision D4 the pollution stock $P$ is the material that has reached the environment and has not regenerated, in gigatonnes, and the CO$_2$ evidence is the bridge that prices it. The loss per gigatonne of CO$_2$ is the marginal loss of the DICE-2023 damage function of \textcite{barrage2024} at the stock the AR6 transient response \parencite{canadell2021} puts at three degrees, the point at which that function is calibrated; block B5 records why the marginal match rather than the level match is the one the pollution price turns on, and the alternative anchor at today's stock is shown beside it. The bridge is the CO$_2$ released per gigatonne of $\Xi$ at the 2000 to 2015 composition of the outflow: the fossil share of domestic processed output in \textcite{haas2020} times the CO$_2$ mass per tonne of fuel at the fuel mix of the same years, from the Global Carbon Budget \parencite{friedlingstein2026} over the fuel masses of block B3. $\kappa$ is the product, with the range over the transient response. Table~\ref{tab:q:data:damages} also gives the output loss at the 1900 stock and the scale the pollution costate should show at 1900, so that a solved path can be checked for a pollution price that has become numerical noise.

Decay is constant: $\theta_0=\theta_{\min}$ and $\theta_P=0$, since block B5 finds the state dependence of the model's form not identified by the impulse response of \textcite{joos2013}, and the sufficient no-tipping condition of the workhorse appendix is then vacuous. The rate is the one that drives the observed emission path to the impulse-response stock in 2024, with the pulse fit as the range's lower end; the two disagree because the driven fit is dominated by recent emissions that have not yet met the slow sinks. $P_0$ is the atmospheric excess stock in 1900 of block B5 converted at the 1900 composition, which is coal and a small fossil share of the outflow. The utility channel is off: DICE-2023 carries no non-market loss separate from its output loss, so a positive $\psi$ would count the same damages twice, and $\varphi$ is then undefined and not written. Land-use CO$_2$ is outside the model.

%% GENERATED by data/build/appendix_tables.py on 2026-09-17
%% Do not edit by hand: the file is rewritten by the next run.
\begin{table}[!htb]
\centering
\begin{threeparttable}
\caption{Damages and decay.}\label{tab:q:data:damages}
\begin{tabular}{lrl}
\toprule
Quantity & Value & Unit \\
\midrule
\multicolumn{3}{l}{\emph{The loss per gigatonne of CO$_2$, DICE-2023 at the 3\,$^\circ$C reference stock}} \\
TCRE best estimate & $9.598\times10^{-6}$ & per GtCO2 \\
TCRE low & $3.386\times10^{-6}$ & per GtCO2 \\
TCRE high & $1.941\times10^{-5}$ & per GtCO2 \\
at today's stock instead (not used) & $2.586\times10^{-6}$ & per GtCO2 \\
\multicolumn{3}{l}{\emph{The bridge from material to CO$_2$}} \\
CO$_2$ per tonne of fuel, 2000--2015 & 2.708 & tCO2/t \\
fossil share of the outflow, 2000--2015 & 0.2608 & share \\
GtCO$_2$ per Gt of $\Xi$ & 0.7062 & GtCO2/Gt \\
$\kappa$ (the point) & $6.777\times10^{-6}$ & per Gt of P \\
$\kappa$, low & $2.391\times10^{-6}$ & per Gt of P \\
$\kappa$, high & $1.370\times10^{-5}$ & per Gt of P \\
\multicolumn{3}{l}{\emph{The 1900 stock}} \\
atmospheric excess stock, 1900 & 27.54 & GtCO2 \\
CO$_2$ per tonne of fuel, 1900 & 2.46 & tCO2/t \\
fossil share of the outflow, 1900 & 0.1418 & share \\
$P_0$ & 78.96 & Gt of material \\
output loss at $P_0$ & $5.350\times10^{-4}$ & share of output \\
$\kappa Y_0/(r+\theta)$, the price scale at 1900 & $2.549\times10^{-4}$ & trillion \$/Gt \\
\multicolumn{3}{l}{\emph{Decay}} \\
$\theta$, emissions-driven fit (the point) & 0.01816 & per year \\
$\theta$, pulse fit (low) & 0.006523 & per year \\
\bottomrule
\end{tabular}
\begin{tablenotes}[flushleft]
\footnotesize
\item[] \textit{Note:} Block B5: Barrage and Nordhaus (2024) with the AR6 transient response, the Global Carbon Budget CO$_2$ by fuel over the B3 fuel masses, and Joos et al.\ (2013). The outflow composition is Haas et al.\ (2020). $\theta_0=\theta_{\min}$ and $\theta_P=0$: the state dependence of decay is not identified. $\psi=0$, so $\varphi$ is undefined and not written.
\end{tablenotes}
\end{threeparttable}
\end{table}


\subsection{Initial stocks}\label{sec:q:data:states}

Table~\ref{tab:q:data:states} collects the six stocks at 1900 with the method each rests on. $K_0$ is the perpetual-inventory stock of block B2, reconstructed and anchored at 1950; $M^K_0$ is the in-use stock of manufactured capital of \textcite{krausmann2018}, the one observed stock; $S_0$ and $X_0$ are the composed reserve of Section~\ref{sec:q:data:extraction}; $P_0$ is the converted stock of Section~\ref{sec:q:data:damages}; and $\mathcal W_0$ is the stockpile consistent with the 1900 disposal flow at the calibrated handling share, since cumulative disposal begins at 1900 by construction and the pre-1900 stockpile is not observed, with the 1900 residue alone as the low end.

%% GENERATED by data/build/appendix_tables.py on 2026-09-17
%% Do not edit by hand: the file is rewritten by the next run.
\begin{table}[!htb]
\centering
\begin{threeparttable}
\caption{The initial stocks at 1900.}\label{tab:q:data:states}
\begin{tabular}{lrrrl}
\toprule
State & Value & Low & High & Method \\
\midrule
$K_0$ & 10.63 & 10.63 & 10.63 & reconstructed (perpetual inventory) \\
$S_0$ & 13,429 & 10,734 & 16,123 & composed, C3 \\
$X_0$ & 13,503 & 10,808 & 16,198 & composed, lower bound \\
$P_0$ & 78.96 & 78.96 & 78.96 & converted, C5 \\
$M^K_0$ & 35.61 & 35.61 & 35.61 & observed \\
$\mathcal W_0$ & 16.39 & 0.572 & 16.39 & reconstructed (flow-consistent) \\
\bottomrule
\end{tabular}
\begin{tablenotes}[flushleft]
\footnotesize
\item[] \textit{Note:} $K_0$ in trillion 2011 international dollars, the rest in gigatonnes; a point where low and high equal the value. Sources are the subsections that identify each stock.
\end{tablenotes}
\end{threeparttable}
\end{table}


\subsection{The metals-only bound}\label{sec:q:data:metals}

Decision D1 reports a metals-only calibration as the bound on the aggregation of Section~\ref{sec:q:cal:aggregation}. It keeps every parameter of the baseline and replaces what the metals evidence identifies on its own: the ceiling at the metals recycling rate, the stock-effect elasticity at the metals price channel, and the reserve, the discovery ceiling and their reference points on the gross-ore basis without the exhaustible-share scaling, since here the one material is metal ore. The initial states that are flows of the whole economy and the extraction level, which is a share of GDP, stay at the baseline: the bound is a bound on the recycling and reserve side, not a metals-only economy. Table~\ref{tab:q:data:metals} gives the entries that differ; the low end of the reserve carries the gap of about half between the two available measurements of metal-ore extraction.

%% GENERATED by data/build/appendix_tables.py on 2026-09-18
%% Do not edit by hand: the file is rewritten by the next run.
\begin{table}[!htb]
\centering
\begin{threeparttable}
\caption{The metals-only bound: entries that differ from the baseline.}\label{tab:q:data:metals}
\begin{tabular}{lrrr}
\toprule
Parameter & Value & Low & High \\
\midrule
$\bar a$ & 0.708 & 0.5269 & 0.8892 \\
$\mu_N$ & 0 & 0 & 1.844 \\
$S_{\mathrm{ref}}$ & 732.1 & 732.1 & 732.1 \\
$X_{\mathrm{ref}}$ & 732.1 & 732.1 & 732.1 \\
$\bar X_{\max}$ & 2,212 & 2,190 & 2,235 \\
$S_0$ & 732.1 & 494.6 & 732.1 \\
$X_0$ & 732.1 & 494.6 & 732.1 \\
\bottomrule
\end{tabular}
\begin{tablenotes}[flushleft]
\footnotesize
\item[] \textit{Note:} \texttt{data/processed/calibration\_metals.json}; every other entry is the baseline's. Metal ores on the gross-ore basis without the exhaustible-share scaling; the low end of the reserve carries the gap between the Krausmann and IRP measurements of metal-ore extraction.
\end{tablenotes}
\end{threeparttable}
\end{table}


\subsection{Checks on the calibrated set}\label{sec:q:data:checks}

The calibrated set is run through the model's own checks from \texttt{model/}, with the calibration file as the only input; the script is \texttt{data/build/c7\_checks.jl} and its output is kept beside the file.

\smalltitle{Maintained restrictions.} \texttt{check\_params} returns no warning: every restriction the theory maintains holds, including the well-posedness condition at the circular rate and the collapse ranking condition. The balanced-dematerialization rates of \theory{Section}{Balanced dematerialization: rate matching} exist, since $\left(\mu_N-1\right)\left(1-\beta_K\right)+\gamma>0$ at the point; they do not exist on the metals-only bound, where $\mu_N$ is at its floor of zero, so that bound has no state B.

\smalltitle{Convexity.} \texttt{convexity\_report} finds four of the conditions of the sufficiency appendix failing, the same four as on the illustrative set: multiplicative damages, the two cost functions, and the intensity condition. The cost functions fail twice over, through the linear term that gives them a finite choke and through stock elasticities above the flow curvature, $\chi_N<\mu_N$ and $\chi_D<\mu_D$, both of which the data put there. Sufficiency is therefore conditional, and the deviation tests of the sufficiency appendix have to be run on every computed path.

\smalltitle{Well-posedness.} \texttt{wellposed\_report} finds all four long-run conditions holding at the circular rate and at the balanced-dematerialization rates, with the cake-eating margin the narrowest.

\smalltitle{The pollution price.} At the calibrated $\kappa$ the output loss at $P_0$ and the price scale $\kappa Y_0/\left(r+\theta\right)$ are those of Table~\ref{tab:q:data:damages}: small, but several orders of magnitude above the Newton tolerance, so the pollution costate at 1900 is a number and not noise.

\smalltitle{A first solve.} The planner corner solved at a short horizon from a fresh starting path reaches a residual four orders of magnitude below its starting value and then stalls on the last steps of the complementarity smoothing, with the treated share at its lower corner in every period: at the municipal collection charge applied to the whole handled flow, no tonne is worth treating, and the recycling margin shuts with it. The same set in the illustrative units stalls at the same place, so it is not a matter of scale. With the collection and treatment charges at a tenth the same solve converges to machine precision, the treated share is interior and rising, and the gate fee changes sign within the horizon. That is a calibration finding as much as a numerical one, and it is the first thing the simulation phase has to settle: the level of the handling charges, which block B4 records as the least transferable number in the waste block, decides whether the loop opens at all.
